% ISSTA 2027 research paper draft. Build notes and rules: README.md in this folder.
\documentclass[sigconf,review]{acmart}
% For a double-blind submission (verify ISSTA 2027 rules): \documentclass[sigconf,review,anonymous]{acmart}

\usepackage{booktabs}
\usepackage{xcolor}
\usepackage{xspace}
\usepackage{listings}

% ---- draft macros -------------------------------------------------------------------------------
\newif\ifanon \anonfalse
\ifanon \newcommand{\tool}{\textsc{Tool}\xspace} \else \newcommand{\tool}{TraceLens\xspace} \fi
\newcommand{\needmacro}[2]{\textcolor{orange}{#2}\typeout{NEEDMACRO: #1 = #2}}
\newcommand{\pending}[1]{\textcolor{blue}{[PENDING: #1]}\typeout{PENDING: #1}}
\newcommand{\fixme}[1]{\textcolor{red}{[FIXME: #1]}\typeout{FIXME: #1}}
\newcommand{\qnx}[1]{\texttt{#1}}
% Generated by research/make_numbers.py. Do not edit.
\newcommand{\tlAACaptureFiveFloorsOff}{4.5\%}
\newcommand{\tlAACaptureTenFloorsOff}{3.2\%}
\newcommand{\tlAAJobFiveFloorsOff}{30.1\%}
\newcommand{\tlAAJobTenFloorsOff}{28.8\%}
\newcommand{\tlBenchBlindAttempted}{3}
\newcommand{\tlBenchBlindReproduced}{0}
\newcommand{\tlBenchCandidate}{7}
\newcommand{\tlBenchExcluded}{6}
\newcommand{\tlBenchReproduced}{3}
\newcommand{\tlBenchTotal}{16}
\newcommand{\tlBusyMediumNoPi}{9 of 9}
\newcommand{\tlBusyMediumPi}{0 of 6}
\newcommand{\tlFamilyBh}{4.3\%}
\newcommand{\tlFamilyHolm}{4.2\%}
\newcommand{\tlFamilyNone}{37.6\%}
\newcommand{\tlGateCapturesPerArm}{30}
\newcommand{\tlGateSplits}{1000}
\newcommand{\tlLinesOfCode}{11265}
\newcommand{\tlMemcachedAfterRate}{0\%}
\newcommand{\tlMemcachedBeforeRate}{100\%}
\newcommand{\tlMemcachedLossySkipped}{1}
\newcommand{\tlMemcachedResidualUpper}{26\%}
\newcommand{\tlMemcachedRunsAfter}{10}
\newcommand{\tlMemcachedRunsBefore}{10}
\newcommand{\tlMemcachedVerdict}{VERIFIED}
\newcommand{\tlMutantsBad}{14}
\newcommand{\tlMutantsGood}{7}
\newcommand{\tlNngAfterRate}{0\%}
\newcommand{\tlNngBeforeRate}{100\%}
\newcommand{\tlNngLossySkipped}{3}
\newcommand{\tlNngResidualUpper}{35\%}
\newcommand{\tlNngRunsAfter}{7}
\newcommand{\tlNngRunsBefore}{10}
\newcommand{\tlNngVerdict}{VERIFIED}
\newcommand{\tlPowerOnePercentFloorsOff}{100\%}
\newcommand{\tlPowerTenPercentFloorsOn}{0\%}
\newcommand{\tlPowerTwentyPercentFloorsOn}{100\%}
\newcommand{\tlRadarHealthy}{0 of 148}
\newcommand{\tlResidualOneOfTen}{39.4\%}
\newcommand{\tlResidualZeroOfFive}{45.1\%}
\newcommand{\tlResidualZeroOfTen}{25.9\%}
\newcommand{\tlResidualZeroOfTwenty}{13.9\%}
\newcommand{\tlRtTestsFindings}{0}
\newcommand{\tlRtTestsRadarFlags}{4}
\newcommand{\tlRtTestsRuns}{16}
\newcommand{\tlRunsNeededTenPercent}{29}
\newcommand{\tlSimPNinetyNineSpreadPct}{0.8\%}
\newcommand{\tlTtkCaught}{100\%}
\newcommand{\tlTtkP}{5.4e-6}
\newcommand{\tlTtkPNinetyNineAfter}{702}
\newcommand{\tlTtkPNinetyNineBefore}{1103}
\newcommand{\tlTtkVerdict}{VERIFIED}
\newcommand{\tlZFourPNinetyNineAfter}{2.99}
\newcommand{\tlZFourPNinetyNineBefore}{7.45}
\newcommand{\tlZOnePNinetyNineAfter}{1.29}
\newcommand{\tlZOnePNinetyNineBefore}{5.74}
\newcommand{\tlZTwoPNinetyNineAfter}{0.77}
\newcommand{\tlZTwoPNinetyNineBefore}{6.75}
\newcommand{\tlZooBoth}{384}
\newcommand{\tlZooBuggyWindows}{883}
\newcommand{\tlZooDetectPerBuild}{10 of 10}
\newcommand{\tlZooFixVerified}{4}
\newcommand{\tlZooFixWorkloads}{5}
\newcommand{\tlZooHealthyRuns}{0 of 30}
\newcommand{\tlZooNeither}{157}
\newcommand{\tlZooRadarOnly}{326}
\newcommand{\tlZooReferenceFlagged}{0 of 50}
\newcommand{\tlZooRulesOnly}{16}
   % generated by research/make_numbers.py; never edit numbers by hand

\lstset{basicstyle=\ttfamily\footnotesize,columns=fullflexible,keepspaces=true}

\setcopyright{none}
\acmConference[ISSTA '27]{ACM SIGSOFT International Symposium on Software Testing and Analysis}{September 2027}{Singapore}
\acmYear{2027}

\begin{document}

\title{Proven on the Kernel: Trace-Validated Repair of Scheduling Bugs on a Microkernel RTOS}

\ifanon
\author{Anonymous Author(s)}
\else
\author{\fixme{authors and affiliations: the user's decision}}
\fi

\begin{abstract}
Priority inversions, IPC deadlocks, lost wake-ups and interrupt-handler overruns in real-time software are found today by an
engineer reading a kernel trace, and a fix is trusted when one re-run looks better. Neither step scales, and the second is
unsound: a scheduling bug that shows in a fraction of runs disappears from a single re-run by chance, and a patch that only
masks a symptom passes a stopwatch. We present \tool, a closed loop for the QNX Neutrino microkernel RTOS that
(i)~proves the root cause of a scheduling bug from kernel events, citing the events that prove it;
(ii)~applies an OS-specific source transform for that root cause only when the transform's safety preconditions hold in the
code, and otherwise gives exact guidance; and
(iii)~accepts the patch only through a statistical, loss-aware oracle that rebuilds the program, re-traces it on the real kernel
over several captures, and states how often the bug could still occur.
\pending{RQ1 headline: false-accept rate of the oracle versus single-run and deadline-only oracles over N mutants}.
On programs we did not write, \tool fixed and verified a semaphore priority inversion in a real-time course exercise and the
lost wake-up of memcached issue \#1348 on QNX~8.0, and it missed an NNG shutdown hang until a rule was added, which we
report as a miss. \pending{RQ2 blind benchmark result, even if null}. Its regression gate, which treats the capture rather than
the job as the statistical unit, holds simulated A/A false alarms at \tlAACaptureFiveFloorsOff{} where pooling jobs alarms in
\tlAAJobFiveFloorsOff{} of builds.
\end{abstract}

\keywords{program repair, real-time systems, kernel tracing, test oracles, priority inversion, RTOS, QNX}

\maketitle

% =================================================================================================
\section{Introduction}
\label{sec:intro}

Real-time systems fail by being late. A control loop that misses its deadline because a lower-priority thread holds the lock it
needs, while a medium-priority thread holds that thread's CPU, has a \emph{priority inversion}; a pair of servers that each send
a message to the other has a \emph{Send/Reply deadlock}; an interrupt thread that does its work before re-enabling the interrupt
delays every later interrupt on that line. These bugs are not visible in the source as a wrong value. They are visible as an
ordering of scheduling events, and only on the target, under its scheduler.

In practice they are found by reading a timeline. The QNX Momentics IDE, for instance, documents how to find a priority
inversion by hand in the System Profiler: inspect the thread states, notice that a higher-priority thread is blocked on a
semaphore while an unrelated thread runs~\cite{qnx-detect-inversion}. Research tools for Linux automate parts of this
diagnosis from LTTng kernel traces~\cite{cote2016problem,desfossez2016latency,tracecompass}, but stop at the diagnosis.
Once an engineer has a fix, it is accepted when a re-run looks better.

Automated repair of concurrency bugs is a mature field. AFix and CFix construct patches for atomicity and order violations from
detector reports~\cite{afix,cfix}; later work fixes races in interrupt-driven embedded code~\cite{sdracer,esfix} and, recently,
uses language models~\cite{drfix}. These systems repair \emph{correctness} bugs and judge a patch by re-running a detector, a
test or a static check~\cite{monperrus2018bibliography}. None that we found repairs a \emph{scheduling} bug, and none uses the
kernel's own record of the patched program's timing, over repeated runs, as the acceptance oracle.
\fixme{the gap claim stays open until commercial RTOS trace tools (Tracealyzer, SystemView, Lauterbach, Green Hills) and
rt-tests literature are searched; see RELATED\_WORK.md search log}

Scheduling bugs make the oracle the hard part. They are intermittent: on our QNX target, planted bugs that depended on
sub-tick release offsets reproduced in one to four runs out of five. They are entangled with platform noise: a virtualised
timer fires a millisecond late, and a host under load turns a 0.13\,ms job into 1.3\,ms. And they can be masked: lowering the
load on the medium-priority thread makes the deadline pass while the inversion remains. A single re-run therefore accepts patches
that do not fix the bug, and a deadline check accepts patches that hide it.

\paragraph{Approach.} \tool closes the loop on the target. It reads a QNX kernel event trace and runs deterministic detectors
that encode QNX semantics (ownerless semaphores have no priority inheritance; message-passing servers inherit their clients'
priority unless the channel is created with \qnx{\_NTO\_CHF\_FIXED\_PRIORITY}; an interrupt stays masked until
\qnx{InterruptUnmask()}). Each finding cites the kernel events that prove it. For each proven root cause a rule rewrites the
source, but only after checking in the code the preconditions under which the rewrite preserves behaviour; otherwise it emits
guidance with source locations. The patched program is rebuilt with the QNX compiler, deployed, and traced again over several
captures. The patch is accepted only if every targeted finding that reproduced before is absent after, no new kind of finding
appears, declared deadlines are met and improved, and lossy captures are excluded rather than scored. The verdict states an
exact upper bound on how often the bug could still occur. If the re-trace shows a next cause, the loop repairs again.

\paragraph{Contributions.}
\begin{itemize}
  \item A closed-loop approach to repairing scheduling bugs in which the kernel trace both proves the cause and decides the fix,
        with precondition-guarded, OS-specific transforms (Section~\ref{sec:approach}).
  \item A statistical, loss-aware acceptance oracle with stated residual-rate bounds, and an evaluation of that oracle by patch
        mutation against single-run and deadline-only oracles (RQ1, Section~\ref{sec:rq1}).
  \item An evaluation on code we did not write, with the tool frozen before the bugs were examined, reporting misses as well as
        fixes (RQ2), and an ablation of the safety preconditions (RQ3).
  \item Evidence that the statistical unit matters: a regression gate that treats a capture as one sample holds its false-alarm
        rate, and one that pools jobs does not (RQ4).
  \item An implementation for QNX~8.0 using only the Python standard library, with its traces, scripts and the numbers in this
        paper regenerated by one command \fixme{artifact link: anonymised mirror if double-blind}.
\end{itemize}

% =================================================================================================
\section{Background: scheduling on QNX Neutrino}
\label{sec:background}

QNX Neutrino is a microkernel RTOS: drivers, file systems and network stacks are user processes that communicate by
synchronous message passing. Its behaviour differs from Linux in ways that decide whether a scheduling pattern is a bug.

\paragraph{Fixed-priority preemptive scheduling.} Threads run at priorities 1--255 under FIFO or round-robin within a priority.
Timers fire on clock-tick boundaries; on QNX~8.0 the tick is fixed at 1\,ms, so two releases less than a tick apart wake in an
order that varies between runs.

\paragraph{Priority inheritance.} QNX mutexes inherit priority by default, so the textbook inversion through a mutex does not
occur. A POSIX semaphore has no owner and therefore no inheritance: a semaphore used as a lock reintroduces unbounded inversion.
Message passing also inherits: a server receiving a message runs at its client's priority, unless its channel was created with
\qnx{\_NTO\_CHF\_FIXED\_PRIORITY}~\cite{qnx-sysarch}.

\paragraph{Message passing.} \qnx{MsgSend()} blocks the client (SEND-blocked, then REPLY-blocked) until the server calls
\qnx{MsgReply()}. Two servers that send to each other while each is the other's only receiver deadlock.

\paragraph{Interrupts.} With \qnx{InterruptAttachEvent()}, an interrupt is delivered to a thread as an event, and the line stays
masked until the thread calls \qnx{InterruptUnmask()}. Work done before unmasking delays every later interrupt on the line.

\paragraph{Tracing.} The instrumented kernel \qnx{procnto-smp-instr} records thread state changes, priority changes,
communication events (send, receive, reply with their receive identifiers), interrupt entry and delivery, and user events
inserted by the program, each with a 64-bit cycle timestamp and CPU. \qnx{tracelogger} writes them to a file and
\qnx{traceprinter} renders them as text.

% =================================================================================================
\section{Approach}
\label{sec:approach}

Figure~\ref{fig:loop} \fixme{draw the loop figure} shows the loop. Its input is a program's source directory, a build
configuration and, optionally, declared deadlines; its output is a patch, an evidence bundle and a verdict.

\subsection{From events to a scheduling model}

The ingest step parses \qnx{traceprinter} text into a normalised event stream, unwrapping the 32-bit cycle counter with the
trace's control-time events and rejecting a trace that does not declare its clock rate rather than guessing it. The model
reconstructs, per thread, a state timeline (running, ready, blocked on mutex, semaphore, condition variable, send, reply,
receive, interrupt wait), a priority timeline from the kernel's own priority-change events (so an inheritance boost is visible
as such), the sync objects each thread waits on, and the client--server pairing of every message from its receive identifier,
which resolves connections made before the capture started.

\paragraph{Lossy traces get no verdict.} A kernel buffer overrun drops events, and a dropped event fakes a long wait. The model
checks scheduler invariants (a thread running on two CPUs at once, a CPU switching threads without the old one leaving) and
reports the violation rate; above a threshold the trace is declared lossy and every downstream consumer treats it as
\emph{no evidence} rather than as clean.

\subsection{Detectors that cite their evidence}

Each detector is a deterministic rule over the model; each finding carries the identifiers of the events that prove it,
the threads and objects involved, and metrics. Severity requires a pattern (the finding repeats, or its magnitude is
large); rare small outliers, which virtualised timing produces, are reported at a lower severity that does not fail a build.

\begin{itemize}
\item \emph{Priority inversion.} A thread $H$ is blocked on a semaphore, a fixed-priority channel or a condition variable whose
  releaser $L$ has lower priority, while $L$ is ready and some $M$ with $p_L < p_M < p_H$ holds $L$'s CPU. Only time during which
  $M$ actually held $L$'s CPU counts. A semaphore used as a lock (waited and posted by the same threads) is told apart from one
  used as a signal.
\item \emph{Send/Reply deadlock and long chains.} A wait-for graph over send- and reply-blocked threads; a cycle is a deadlock unless
  some thread of a process in the cycle is receiving on the very channel being called.
\item \emph{Mutex contention.} An uncontended QNX mutex never enters the kernel, so every mutex-blocked interval is contention; waits
  on locks at kernel addresses are reported but never fail a build.
\item \emph{Interrupt handler overrun.} Time from an interrupt thread's return from \qnx{InterruptWait()} to its
  \qnx{InterruptUnmask()}, per line, with shared lines keeping all their handlers.
\item \emph{Deadline miss.} Response time from a job's release, with releases fitted per stream by Theil--Sen
  regression~\cite{sen1968} because the trace clock and the program clock drift apart (0.05\% on our target); each late job is
  broken down into late wake-up, waiting for a CPU, blocked and running, and a late wake-up while nothing of equal or higher
  priority ran is attributed to the platform.
\item \emph{Lost wake-up} and \emph{stuck job}, added during the evaluation (Section~\ref{sec:rq2}) and labelled as such: a
  condition variable signalled with no waiter, followed by a wait still blocked at the end of the trace while another thread joins
  the waiter; and a job started with no end past its budget.
\end{itemize}

\subsection{Precondition-guarded remedies}

A remedy maps one proven root cause to one source transform (Table~\ref{tab:remedies}). It runs on a sandbox copy containing
only the program's own files, parses the relevant C constructs, and applies the transform only if its preconditions hold;
otherwise it returns guidance naming the source lines and the change to make. The preconditions are what make a transform
behaviour-preserving apart from timing; RQ3 measures what happens without them.

\begin{table*}[t]
\caption{Remedies. Each is applied automatically only when its preconditions hold in the source.}
\label{tab:remedies}
\small
\begin{tabular}{@{}p{0.2\textwidth}p{0.36\textwidth}p{0.38\textwidth}@{}}
\toprule
Root cause (from the trace) & Transform & Applied automatically when \\
\midrule
Semaphore used as a lock & \qnx{sem\_t} $\to$ \qnx{pthread\_mutex\_t} with \qnx{PTHREAD\_PRIO\_INHERIT}; wait/post $\to$ lock/unlock &
  initial value 1; every function that waits also posts; no trywait, timedwait or other use \\
Fixed-priority server channel & drop \qnx{\_NTO\_CHF\_FIXED\_PRIORITY}: the server inherits its clients' priority &
  the flag appears in a \qnx{ChannelCreate()} flag expression \\
Work before \qnx{InterruptUnmask()} & unmask first; move the work to a worker fed by a coalescing pulse, at a priority below every
  interrupt thread seen in the trace, started before attach and stopped after detach &
  the work uses no locals of the loop's function and no device access follows it \\
Private work inside a contended lock & move leading/trailing statements that touch no shared data out of the critical section &
  the statements use no file-scope data and no pointer that may alias it; no return, break or goto \\
Lost wake-up at shutdown & re-check the run flag before \qnx{pthread\_cond\_wait()} & the stopper sets the flag and signals
  while holding the same mutex \\
Mutual \qnx{MsgSend()} & guidance only: a send hierarchy, one direction becomes a pulse & never (the reply carries data) \\
\bottomrule
\end{tabular}
\end{table*}

\subsection{The acceptance oracle}
\label{sec:oracle}

The oracle compares $n_b$ captures of the original build with $n_a$ captures of the patched build. Lossy captures are excluded and
counted; with fewer than five clean captures on either side the verdict is INCONCLUSIVE. A finding kind is a \emph{target} if
it appears in at least 60\% of clean before-captures; a target that did not reproduce cannot be verified fixed, which prevents
a fix from being credited for a bug that never happened. The verdict is VERIFIED only if
(1)~every target is absent from every clean after-capture,
(2)~no finding kind appears after that never appeared before, and
(3)~for each declared deadline the after p99 is within budget and, where the before-captures missed it, a one-sided
Mann--Whitney U test~\cite{mannwhitney1947} shows the after response times to be smaller at $\alpha = 0.01$ (exact permutation
distribution up to 60 samples, normal approximation with tie correction above).
Otherwise the verdict is NOT FIXED. \fixme{the Mann--Whitney test in (3) currently pools all jobs of all captures
(verdict.py); RQ4 shows job pooling inflates significance under A/A. Either move (3) to per-capture p99, like the gate, or
justify pooling; R2 compares both (O4-job vs O4-capture)}

\paragraph{Stating the residual rate.} Absence in $n_a$ captures does not prove absence. If a target was seen in $k$ of $n_a$
clean after-captures, the verdict reports the one-sided Clopper--Pearson upper bound~\cite{clopperpearson1934} on its
per-capture rate; for $k = 0$ this is $1-(1-c)^{1/n_a}$. At 95\% confidence, zero sightings in 5, 10 and 20 captures bound the
remaining rate by \tlResidualZeroOfFive, \tlResidualZeroOfTen{} and \tlResidualZeroOfTwenty; bounding it by 10\% takes
\tlRunsNeededTenPercent{} captures. The unit is the capture: periods within one capture share a boot, a host load and a cache
state, so they do not tighten the bound.

\subsection{Repairing again}

A patch can remove the proven cause and expose the next. On one workload the priority-inheriting mutex removed the inversion and
the re-trace then showed the high-priority thread waiting about a millisecond each period on the low-priority thread's long
critical section. On NOT FIXED, \tool diagnoses the patched build's own captures and applies the remedy for the newly proven
cause, for at most three rounds; each round's patch, build log, traces and verdict are kept in the evidence bundle.

\subsection{Regression gate}

The same machinery guards later builds. A baseline records, per declared deadline and per job stream, the p99 of each capture.
A candidate build fails if a finding kind appears that the baseline did not have, if a deadline the baseline met is now missed,
or if a one-sided exact Mann--Whitney test on per-capture p99 shows it slower, with Holm's correction~\cite{holm1979} across
all deadlines and an effect floor (10\% and 50\,\textmu{}s by default). With fewer captures than a rank test can use, the gate
falls back to a stated threshold and reports its exact chance of passing an unchanged build.

% =================================================================================================
\section{Implementation}
\label{sec:impl}

\tool is \tlLinesOfCode{} lines of Python using only the standard library. It manages a QNX~8.0 virtual
machine built with \qnx{mkqnximage}, copies a \qnx{qcc}-built binary to the target, snapshots thread priorities with
\qnx{pidin}, runs the program under \qnx{tracelogger}, and fetches and decodes the trace. Reports are produced as terminal
text, HTML with per-thread timelines, JSON, SARIF and Chrome trace format. \fixme{decide whether to describe the learned
radar here or only in the discussion}

% =================================================================================================
\section{Evaluation}
\label{sec:eval}

\begin{description}
\item[RQ1] How reliable is the oracle: how often does it accept a bad patch or reject a good one, as a function of the number of
  captures, compared with simpler oracles?
\item[RQ2] Do the detectors and remedies work on code we did not write?
\item[RQ3] Do the remedies' preconditions matter?
\item[RQ4] Does treating the capture as the statistical unit control false alarms?
\end{description}

\paragraph{Platform.} QNX SDP~8.0, instrumented kernel \qnx{procnto-smp-instr}, on a QEMU/KVM x86\_64 virtual machine with four
virtual CPUs, hosted under WSL2 on Windows~11; programs built with \qnx{qcc} (gcc~12.2.0). Captures were taken with the host
idle, because host load appears inside the guest as running time. \fixme{hardware results for the revision; emulated aarch64
only as a portability note}

\paragraph{Development set.} Five small QNX programs with planted bugs (Table~\ref{tab:zoo}) were used to build and calibrate the
detectors and remedies. We report them as a development set, not as an evaluation: the bugs and the rules were written by the
same authors. On a held-out ten-capture run, each buggy build was caught in \tlZooDetectPerBuild{} captures,
the hand-written reference fixes were flagged in \tlZooReferenceFlagged{} captures, and three healthy system programs in
\tlZooHealthyRuns{}; \tlZooFixVerified{} of \tlZooFixWorkloads{} were repaired automatically
and VERIFIED, and the deadlock received guidance.

\begin{table}[t]
\caption{Development set (planted bugs; not an evaluation). p99 response before $\to$ after the automatic fix, 5 captures a side.}
\label{tab:zoo}
\small
\begin{tabular}{@{}lll@{}}
\toprule
Program & Planted bug & Before $\to$ after \\
\midrule
Infusion pump & semaphore as lock & \tlZOnePNinetyNineBefore{} $\to$ \tlZOnePNinetyNineAfter{}\,ms \\
Brake-by-wire & fixed-priority channel & \tlZTwoPNinetyNineBefore{} $\to$ \tlZTwoPNinetyNineAfter{}\,ms \\
Sensor fusion & mutual \qnx{MsgSend()} & guidance only \\
Camera pipeline & work inside global lock & \tlZFourPNinetyNineBefore{} $\to$ \tlZFourPNinetyNineAfter{}\,ms \\
Motor encoder & work before unmask & \pending{Z5 masked time, from the z5 bundle} \\
\bottomrule
\end{tabular}
\end{table}

% -------------------------------------------------------------------------------------------------
\subsection{RQ1: Reliability of the oracle}
\label{sec:rq1}

\paragraph{Design.} We wrote patches whose effect we know and asked each oracle to judge them. For each development program,
the hand-written reference fix is a \emph{good} patch. \emph{Bad} patches are of six kinds: no-ops; fixes with the wrong
mechanism (a mutex explicitly without inheritance; a server that re-lowers its own priority after every receive); symptom
masking (less load on the interfering thread, a shorter timeout on a deadlocked send); rarer races (release offsets that make
the inversion appear in about a fifth of captures); fixes that introduce a new cause (a priority-inheriting mutex held across a
slow write); and a real past mistake of our own remedy, which ran deferred interrupt work one priority below the interrupt thread
and started it after attaching, delaying other interrupts by up to 20\,ms. One patch is a functional bug that timing cannot see
(a shared copy moved out of its lock); we expected every timing oracle to accept it. The set also covers the two third-party
programs of RQ2 (their own fixes as good patches, and for memcached a plausible wrong fix that signals twice instead of
re-checking the flag): \needmacro{MutantsGood}{7} good and \needmacro{MutantsBad}{14} bad patches in all. For each, the expected
verdict of every oracle, with a stated confidence, was written and frozen before any capture (protocol commit 8f07f01).
Each patch was captured 20 times.
We compare six oracles: O1 a single clean re-run; O2 after-p99 within budget; O3 targets absent in all captures, no statistics;
O4 \tool's oracle as shipped, whose deadline test pools the jobs of all captures; O4c the same oracle with its deadline test on
one p99 per capture; O5 O4 with the residual bound required below 20\%. For $N \in \{1,3,5,10,20\}$ we draw 1{,}000 subsets of
$N$ captures per side.

\paragraph{Results.} \pending{Table: false-accept and false-reject rate per oracle and N; figure: false-accept versus N. Expected
(frozen): O1 and O2 accept the masking and rare-race patches; O4 rejects them; the rare-race false-accept rate of O3/O4 falls with
N; every oracle accepts the data-race patch.}

% -------------------------------------------------------------------------------------------------
\subsection{RQ2: Code we did not write}
\label{sec:rq2}

We distinguish \emph{case studies}, in which we examined the bug before running the tool or changed the tool afterwards, from a
\emph{blind} evaluation, in which the tool was frozen (git tag and source hash recorded) before any bug was examined and each
bug's driver shows the symptom by its own exit code.

\paragraph{Case studies.}
\emph{A real-time course exercise.} A teaching program from a university real-time systems course runs three
\qnx{SCHED\_RR} threads around a global semaphore used as a lock; the course asks students to observe the inversion. Unedited
apart from a header-level port layer, \tool reported a critical priority inversion in \tlTtkCaught{}
of captures, applied the semaphore-to-mutex remedy in one round, and the fix was VERIFIED: the inversion was absent in every after-capture
and the high-priority job's p99 fell from \tlTtkPNinetyNineBefore{} to \tlTtkPNinetyNineAfter{}\,ms
(one-sided exact Mann--Whitney $p = \tlTtkP{}$), with the configuration frozen before the confirmation
run. The fix is the exercise's intended answer.

\emph{memcached \#1348.} \qnx{stop\_logger\_thread()} sets a flag and signals a condition variable; if the signal arrives while
the logger thread is in its last sleep, the logger later waits forever and shutdown hangs. With memcached's sources unedited and
built for QNX~8.0 (the whole daemon with libevent, and separately the logger alone), the hang reproduced in
\tlMemcachedBeforeRate{} of captures and the one-line fix (re-check the flag before waiting) was \tlMemcachedVerdict:
\tlMemcachedAfterRate{} of \tlMemcachedRunsAfter{} after-captures, residual bound \tlMemcachedResidualUpper. The detector and the
remedy for this pattern were written \emph{after} we had seen the bug; this demonstrates the loop on production code, not a
detection rate.

\emph{NNG \qnx{nng\_close()}.} On a long-discussed shutdown hang in the NNG messaging library, built for QNX from the release
before the upstream fix and from the fix commit, the tool as it stood reported nothing: a miss. A stuck-job rule added afterwards
reports the hang in \tlNngBeforeRate{} of buggy captures and \tlNngAfterRate{} of the fixed one's
(\tlNngRunsAfter{} clean after-captures, \tlNngLossySkipped{} lossy excluded, residual bound \tlNngResidualUpper). The repair is
upstream's, not ours. \pending{preregistered residual-hang check on the fixed commit (docs/results/nng\_close/RESIDUAL.md)}

\emph{A clean result.} Three \qnx{rt-tests} programs (\qnx{pi\_stress}, \qnx{ptsematest}, \qnx{pmqtest}) built behind a
compatibility layer produced \tlRtTestsFindings{} rule findings in \tlRtTestsRuns{} runs (the learned detector of
Section~\ref{sec:discussion} raised \tlRtTestsRadarFlags{} medium-severity flags on \qnx{ptsematest}), while a control that
sustains an inversion on purpose was flagged in \tlBusyMediumNoPi{} runs without priority inheritance and
\tlBusyMediumPi{} with it.

\paragraph{Blind evaluation.} Of \tlBenchTotal{} bugs considered so far, \tlBenchReproduced{} are the case studies above,
\tlBenchExcluded{} were excluded (most often because the bug did not reproduce on QNX; a preregistered attempt on a broker hang
gave no hang in 39 runs), and \tlBenchCandidate{} are candidates. Of \tlBenchBlindAttempted{} attempts made under a freeze,
\tlBenchBlindReproduced{} reproduced. \pending{blind results under research-freeze-1, rubric D1--D3, F0, R1--R5; report misses}

% -------------------------------------------------------------------------------------------------
\subsection{RQ3: Do the preconditions matter?}
\label{sec:rq3}

\pending{R5.1: each remedy applied with its preconditions ignored, on the development set and on a suite of programs in which
the naive transform is wrong (a semaphore used as a signal; wait and post in different functions; a counting semaphore; timed
waits; device access after deferred interrupt work; an aliased pointer to shared data). Metrics: compile failures, functional
test failures, new findings, VERIFIED rate. R5.2: one round versus up to three.}

% -------------------------------------------------------------------------------------------------
\subsection{RQ4: The statistical unit}
\label{sec:rq4}

In the regression gate we compared two ways of testing whether a build is slower: one p99 per capture (the gate), and every job
of every capture pooled. On \tlGateCapturesPerArm{} simulated captures of an unchanged build and \tlGateSplits{} random splits,
the capture-level test raised false alarms in \tlAACaptureFiveFloorsOff{} of 5-versus-5 and \tlAACaptureTenFloorsOff{} of
10-versus-10 comparisons, at $\alpha = 0.05$; the pooled-job test alarmed in \tlAAJobFiveFloorsOff{} and \tlAAJobTenFloorsOff.
The roughly 200 jobs of a capture share its boot, host load and cache state; counting them as independent samples is what makes
the pooled test cry wolf. Across ten deadlines tested at once, uncorrected tests raised at least one false alarm in
\tlFamilyNone{} of comparisons, Holm's correction in \tlFamilyHolm{}. The simulator is quieter than the target (capture-to-capture
p99 spread \tlSimPNinetyNineSpreadPct), so it cannot measure power, and we make no power claim. \pending{VM part of R8: A/A and an
injected slowdown on the target}

% =================================================================================================
\section{Discussion}
\label{sec:discussion}

\paragraph{What real QNX taught the tool.} Several design decisions came from the target, not from the specification: user
events that some \qnx{traceprinter} builds print only as raw payload words; planted bugs that reproduced one to four times in five
until releases were spaced by whole ticks; a rigid release grid that turned 0.05\% clock drift into 127 false misses in 200 jobs;
a deferred-work remedy that had to run below every interrupt thread in the trace, including a virtio driver's; and an oracle that
once verified a fix for a deadlock that had not reproduced. \fixme{trim; the long version is an OSPERT paper}

\paragraph{A learned second opinion did not transfer.} A detector trained only on healthy traces flagged
\tlZooRadarOnly{} of \tlZooBuggyWindows{} held-out buggy captures of bug families no rule covers, with
\tlRadarHealthy{} false alarms, on traces from our own generator. On the first two real programs it flagged the
hung and the healthy run alike. We keep it out of the loop: it never gates a build and never produces a patch.

\paragraph{What the oracle cannot see.} The oracle judges timing and scheduling findings; a patch that removes a lock around
shared data can pass it. This is why the remedies check their preconditions in the source rather than relying on the oracle
(RQ3), and why RQ1 includes such a patch.

% =================================================================================================
\section{Threats to Validity}
\label{sec:threats}

\emph{Construct.} The development set was written by the authors; we report it as development data. Two detectors used in
RQ2 (lost wake-up, stuck job) were written after the bugs they detect were seen, and are labelled as such.
\emph{Internal.} Timing on a virtual machine nested under WSL2 is noisier than on hardware; we capture with the host idle,
exclude lossy captures and require repeated captures, but absolute response times are platform-specific. The instrumented
kernel and the tracer may perturb timing \pending{probe-effect measurement, R4.2}.
\emph{External.} One RTOS, one architecture on a virtual machine; an emulated aarch64 guest runs the same suite but is not
evidence about ARM hardware. \pending{hardware}.
\emph{Conclusion.} With five captures per side the residual bound is loose (\tlResidualZeroOfFive); we report it rather than
claim absence. Overlapping subsamples in RQ1 and RQ4 make the reported intervals somewhat narrow.

% =================================================================================================
\section{Related Work}
\label{sec:related}

\paragraph{Trace analysis for real-time systems.} Côté and Dagenais analyse LTTng kernel traces of periodic real-time tasks
on Linux, combining critical-path analysis with scheduling information~\cite{cote2016problem}; Desfossez et al.\ track latency
in production kernels~\cite{desfossez2016latency}; Trace Compass provides latency and critical-path views~\cite{tracecompass};
Fournier et al.\ combine user and kernel traces to explain slow web requests~\cite{fournier2019cause}. These diagnose; \tool
also repairs and verifies, and encodes QNX semantics that Linux traces do not expose in the same form (ownerless semaphores,
fixed-priority channels, interrupt masking by an interrupt thread).

\paragraph{Concurrency bug repair.} AFix and CFix build patches from detector reports for atomicity and order
violations~\cite{afix,cfix}; RaceFixer starts from ThreadSanitizer reports~\cite{racefixer}; SDRacer and ESfix repair races in
interrupt-driven embedded code~\cite{sdracer,esfix}; Dr.Fix uses a language model with retrieved fixes~\cite{drfix}. They target
correctness bugs, and their patches are judged by re-detection, tests or static checks~\cite{monperrus2018bibliography}. \tool
targets scheduling bugs, and its oracle is the kernel's record of the patched program over repeated captures.

\paragraph{Statistics for performance testing.} Rigorous benchmarking accounts for variation between runs and between
repetitions~\cite{kalibera2013rigorous}; change-point detection finds regressions in performance histories~\cite{daly2020changepoint}.
We adopt the same principle (the run, not the iteration, is the sample) for both the acceptance oracle and the gate, and add
residual-rate bounds for intermittent bugs.

% =================================================================================================
\section{Conclusion}
\label{sec:conclusion}

\pending{write after RQ1--RQ3 results}

\bibliographystyle{ACM-Reference-Format}
\bibliography{refs}

\end{document}
